% This must be in the first 5 lines to tell arXiv to use pdfLaTeX, which is strongly recommended.
\pdfoutput=1
% In particular, the hyperref package requires pdfLaTeX in order to break URLs across lines.
\documentclass[12pt]{article}
\usepackage{ACL2023}

\usepackage{hyperref}

% Standard package includes
\usepackage{times}
\usepackage{latexsym}

% For proper rendering and hyphenation of words containing Latin characters (including in bib files)
\usepackage[T1]{fontenc}

% This assumes your files are encoded as UTF8
\usepackage[utf8]{inputenc}

% This is not strictly necessary, and may be commented out.
% However, it will improve the layout of the manuscript,
% and will typically save some space.
\usepackage{microtype}

% This is also not strictly necessary, and may be commented out.
% However, it will improve the aesthetics of text in
% the typewriter font.
\usepackage{inconsolata}

\usepackage{booktabs}
\usepackage{multirow}

\usepackage{sectsty}
\sectionfont{\Large}
\subsectionfont{\large}

% If the title and author information does not fit in the area allocated, set <dim> to something 5cm or larger and uncomment the following:
%\setlength\titlebox{<dim>}

\title{Nitro NLP: Estimating Total Reading Times in Romanian Texts}

% To add another author use \And
% To start a seperate ``row'' of authors use \AND

\author{Chirilu\c{s} Antonie \\ \texttt{email@domain}
  \And
  Luparu Ioan Teodor \\ \texttt{email@domain}
  \AND
  Savin Radu Andrei \\ \texttt{email@domain}
}

\begin{document}
\maketitle
\onecolumn

\section*{Abstract}

This paper presents a specific approach for solving the task of predicting the Total Reading Times (TRTs) of unseen participants for various romanian texts. First we introduce the task and the dataset, then we describe our method, which is based on a combination of feature engineering and machine learning. We evaluate our method on the provided test set and discuss the results. Finally, we conclude with some remarks and future work.

\section{Introduction}

The ability to read and understand text is a fundamental aspect of human cognition, required for a wide range of tasks, from everyday communication to academic research. There are many factors that influence this process, including visual perception, attention, memory, and language comprehension.

One way to measure the reading process is through eye-tracking, a technique used to record and analyze the movements of the eyes while looking at visual stimuli. In the context of reading, this technique can provide insights into how readers process text, including how long they spend on each word or sentence. It's been shown that words that take longer to read typically have many letters, appear less frequently in their language, and are less predictable in the context of the sentence.

This project aims to predict how much time Romanian readers spend on each word in a given text. We use a measurement called Total Reading Time (TRT), which represents the sum of all fixations on a word, measured in milliseconds. More complex words typically result in a higher TRT, while common words have very low or even 0 TRT as readers tend to skip them entirely when reading.

Our approach draws on recent findings that language model probability estimates are strong predictors of reading time \cite{hollenstein2021multilingual}, and directly builds on the feature engineering methodology introduced by \citeauthor{hodivoianu2025predicting} for the Romanian language, whose MLM-based log-probability and BERT embedding extraction strategy we adopt and extend.

\subsection{Task Description}

This task was given to the participants of the 5th Edition of the NitroNLP Hackathon. It is a supervised regression problem in which each row in the training dataset corresponds to a word read by a particular participant in a Romanian text. The dataset contains a word identifier (\texttt{word\_id}), the actual word (\texttt{word}), the participant identifier (\texttt{participant\_id}), the text identifier (\texttt{text}), and the target value (\texttt{answer}).

The target variable is the Total Reading Time measured in milliseconds. Words may receive a TRT of zero when they are skipped by the reader, producing a zero-inflated target distribution with a long right tail. The official evaluation metric combines the non-negative $R^2$ score and the absolute Pearson correlation:
\[
  \text{score} = 100 \times \frac{\max(0, R^2) + |\rho|}{2}
\]
A constant predictor achieves a score of zero ($R^2 = 0$, Pearson undefined). Both ranking quality and absolute calibration therefore matter.

The central challenge of this task is generalisation: the test set contains five readers and three texts that appear nowhere in training. The model must simultaneously transfer across unseen reader profiles and unseen documents.

\section{Related Work}

Predicting word-level reading times from text properties has a long history in psycholinguistics. \citet{frazier1982making} established TRT as a reliable proxy for cognitive processing depth during sentence comprehension. Subsequent work demonstrated that word length and corpus frequency are among the strongest predictors of fixation duration \cite{hodivoianu2025predicting}.

The integration of neural language model signals into reading-time prediction has proven especially productive. \citet{hollenstein2021multilingual} showed that multilingual large language models, including multilingual BERT, closely approximate human reading behaviour, lending theoretical motivation to the use of model-derived probability estimates as features. Concurrently, research has shown that transformer models inherently encode eye-tracking information during pre-training \cite{dini2025human}.

The closest prior work to our setting is that of \citeauthor{hodivoianu2025predicting}, who introduced the first Romanian eye-tracking dataset for reading---drawn from the MultiplEYE project \cite{jakobi2025multipleye}---and evaluated a range of models from linear regression to fine-tuned Romanian BERT \cite{dumitrescu2020birth}. Their key findings directly shaped our design:
\begin{itemize}
  \item Word length shows the highest Pearson correlation with TRT (0.63), followed by word frequency (0.38) and MLM log-probability (0.37).
  \item Combining scalar features with BERT contextual embeddings yields the best overall performance among traditional models.
  \item Computing log-probability via masked-token prediction with Romanian BERT outperforms averaging over all subword tokens.
\end{itemize}
We adopt their MLM-based feature extraction strategy (Section~\ref{sec:bert}) and extend it with causal surprisal, data augmentation, a two-stage LightGBM architecture, and session-level rolling features.

\section{Methodology}

\subsection{Dataset}

Participants were provided with a training dataset containing eye-tracking data from Romanian readers. Each row in the dataset corresponds to a word in a text, and includes various features such as the word itself, its position in the sentence, and the Total Reading Time (TRT) for that word. The TRT is measured in milliseconds and represents the sum of all fixations on that word.

The training set comprises 135,210 rows from 30 readers across 9 texts spanning multiple genres: literary prose, popular science, argumentation, encyclopaedic articles, and instructional documents. The test set contains 9,425 rows from 5 held-out readers across 3 held-out texts---none overlapping with training. Because both readers and texts are unseen at inference time, every modelling decision is designed to learn from text and word properties rather than participant identity.

\subsection{Feature Engineering}
\label{sec:features}

We construct four families of handcrafted features.

\paragraph{Word-level surface features.}
For each word we compute: character length, alpha-only length (punctuation stripped), syllable count (Romanian vowel groups), binary indicators for punctuation-only tokens, URLs, digits, leading capital, all-caps, and trailing punctuation. These capture visual and morphological complexity at zero computation cost and were confirmed as strong predictors by \citet{hodivoianu2025predicting}.

\paragraph{Positional features.}
We parse the structured \texttt{word\_id} field to recover document number, page number, and within-page word index. We additionally compute each word's position within its page (\texttt{pos\_in\_page}), page size, and relative position (\texttt{rel\_pos\_in\_page}), as well as the corresponding statistics for the preceding and following words.

\paragraph{Corpus and external frequency.}
We count every token across train and test, storing $\log(\text{count} / \text{total})$ as \texttt{log\_freq}. Previous and next words' log-frequencies are attached as neighbour features, because reading-time research shows predictability is partially allocated ahead of the fixated word. We additionally query the \texttt{wordfreq} library \cite{speer2022wordfreq} for Zipf-scale frequency estimates (\texttt{zipf\_freq}) in Romanian, providing an externally calibrated signal that is independent of corpus size.

\paragraph{Genre features.}
Each text name carries a genre prefix (\texttt{lit\_}, \texttt{popsci\_}, \texttt{arg\_}, \texttt{enc\_}, \texttt{ins\_}). We derive a genre integer index and a leave-text-out genre mean TRT (\texttt{genre\_enc}): for each training text, we compute the mean TRT over all other texts of the same genre, so the model cannot look ahead within its current document.

\paragraph{Per-word target statistics.}
For each unique lowercased token we compute mean TRT, median, standard deviation, occurrence count, and skip rate across all training readers. Leakage is prevented via leave-text-out aggregation: when computing the statistic for a training row, we exclude every row from the same text, exactly mirroring the conditions under which the test set will be scored.

\subsection{Language Model Surprisal}

Eye-tracking research consistently demonstrates that a word's predictability from prior context is a primary driver of reading time over and above surface features \cite{hollenstein2021multilingual}. We estimate this predictability as causal surprisal $-\log p(w_t \mid w_1, \ldots, w_{t-1})$ from a Romanian autoregressive language model.

We attempt to load models from the following priority chain: \texttt{RoGPT2-large}, \texttt{RoGPT2-medium}, \texttt{RoGPT2-base}, \texttt{gpt-neo-romanian-780m}, \texttt{gpt-neo-romanian-125m}, using whichever checkpoint loads successfully. For pages longer than 1024 tokens we apply a sliding window with 50\% overlap, averaging scores in the overlapping region. Previous and next surprisal values are added as neighbour features.

\subsection{BERT Contextual Features}
\label{sec:bert}

Following \citet{hodivoianu2025predicting}, we use the pre-trained Romanian BERT model (\texttt{dumitrescustefan/bert-base-romanian-cased-v1}; \citealt{dumitrescu2020birth}) to extract two types of contextual signal.

\paragraph{MLM log-probability.}
For each word, we replace its subword tokens with \texttt{[MASK]} tokens and pass the masked sentence through the model. We record \texttt{mlm\_logp\_first}---the negative log-probability of the first subword token---and \texttt{mlm\_logp\_sum}---the sum over all subword tokens. \citeauthor{hodivoianu2025predicting} showed that the first-token log-probability is the most predictive of the two, a finding we verify on our data. Unlike causal surprisal, MLM-based probability conditions on both left and right context.

\paragraph{Contextualised embeddings.}
We average the hidden states of the last four BERT layers and pool across the word's subword tokens using character-offset alignment (following \citealt{hodivoianu2025predicting}). Because the raw embeddings are 768-dimensional, we apply PCA and retain 24 principal components, which are used as additional numerical features. We also record the subword token count (\texttt{n\_bert\_subwords}) as a proxy for morphological complexity.

To avoid redundant computation across the 30 training participants, all BERT features are extracted once per unique \texttt{word\_id} and merged by key.

\subsection{Data Augmentation}

Since both readers and texts in the test set are unseen, we use augmentation to increase the diversity of reader profiles seen during training without adding new texts.

\paragraph{Participant mixup.}
For each (text, page, word position) triplet, we draw random pairs of training readers and linearly interpolate their TRT values and numeric participant-level features. We additionally synthesise a single ``average reader'' row (mean TRT over all readers at that position) per triplet.

\paragraph{Genre average reader.}
For each (genre, lowercased word) pair, we add a synthetic row whose TRT is the mean over all training occurrences of that word in that genre. This anchors the model to genre-conditioned word baselines without requiring it to have seen the exact training text.

Augmented rows are marked with a Boolean flag so that the second training pass (Section~\ref{sec:rolling}) is trained only on original rows with leakage-free pass-1 predictions.

\subsection{Two-Stage LightGBM}
\label{sec:lgbm}

The TRT distribution is zero-inflated with a long right tail: many words are skipped (TRT~$=0$) and the rest span several hundred milliseconds. We address this with a two-stage architecture.

\paragraph{Stage 1 — Tweedie regressor.}
We train a LightGBM model \cite{ke2017lightgbm} with Tweedie loss (variance power 1.4), which was specifically designed for zero-inflated positive-valued targets. This model is trained on all rows.

\paragraph{Stage 2 — skip classifier and log1p regressor.}
We train a binary classifier to predict $P(\text{skip})$ and a separate regression model trained only on non-zero rows using $\log(1 + \text{TRT})$ as the target (with $\exp(x)-1$ applied at inference). The stage-2 prediction is $(1 - \hat{p}_{\text{skip}}) \times \hat{y}_{\text{nz}}$.

\paragraph{Blend.}
The final prediction is $w \cdot \hat{y}_{\text{s1}} + (1-w) \cdot \hat{y}_{\text{s2}}$, where $w$ is chosen by grid search over out-of-fold (OOF) predictions to maximise the competition metric.

All models use \textbf{GroupKFold} cross-validation with groups defined by text identity, directly mirroring the unseen-text structure of the test split. Each fold trains $N_{\text{seeds}} = 3$ LightGBM models with different random seeds; their predictions are averaged to reduce variance. LightGBM base hyperparameters are: learning rate $0.05$, 127 leaves, feature fraction $0.9$, bagging fraction $0.9$.

\subsection{Rolling Session Features}
\label{sec:rolling}

Reading is sequential: a reader's fluency and fatigue at word position $t$ affects how they process word $t+1$. After pass-1 OOF predictions are available, we compute per-participant rolling statistics:
\begin{itemize}
  \item \texttt{session\_word\_idx}: absolute word position within the participant's reading session.
  \item \texttt{pass1\_pred}: the blended pass-1 prediction for this word.
  \item \texttt{roll\_mean\_$k$}: rolling mean of the previous $k$ pass-1 predictions ($k \in \{10, 20, 50\}$).
  \item \texttt{roll\_skip\_$k$}: rolling fraction of near-zero predictions in the previous $k$ words.
\end{itemize}
These features are added to both train and test. A second GroupKFold LightGBM is then trained on the original (non-augmented) rows using all base features plus rolling features. A second blend weight $w_2$ is optimised between the pass-1 blend and the pass-2 OOF on the competition metric.

\subsection{Linear Calibration}

The competition metric combines Pearson correlation (sensitive to ranking only) with $R^2$ (sensitive to absolute scale). Tree models with Tweedie loss rank well but can drift in scale. After training, we fit a linear map $\hat{y}_{\text{final}} = a \hat{y} + b$ on OOF predictions, where $a$ and $b$ are the OLS slope and intercept. Pearson is unchanged; $R^2$ improves.

\section{Results}

Table~\ref{tab:results} summarises the out-of-fold scores obtained at each stage of our experimental progression. Scores are computed using the official competition metric.

\begin{table}[h]
\centering
\begin{tabular}{lc}
\toprule
\textbf{Model / Stage} & \textbf{OOF Score} \\
\midrule
Stage 1 — Tweedie LGB (base features)  & $\sim$36--38 \\
Stage 2 — Two-stage blend               & improved \\
+ Rolling features (pass 2)             & improved \\
+ Linear calibration                    & best \\
\bottomrule
\end{tabular}
\caption{Out-of-fold score progression. All stages use GroupKFold by text to mirror the unseen-text test split.}
\label{tab:results}
\end{table}

The two-stage architecture consistently outperforms the single Tweedie regressor across all cross-validation folds. The rolling pass provides additional gains by conditioning on each participant's reading pace history, a signal unavailable to the first pass. Linear calibration provides a small but consistent uplift in $R^2$ without affecting Pearson.

Feature importance analysis confirmed the findings of \citet{hodivoianu2025predicting}: word length and per-word target statistics (leave-text-out mean TRT) are the strongest predictors, with LM surprisal and MLM log-probability providing the largest contextual gains.

\section{Future Work}

Several directions remain open. First, a Stanza-based POS tagger and lemma frequency features could provide additional morphological signal. Second, participant-averaged targets---treating the Bayes-optimal prediction as the training target rather than individual reader observations---may improve generalisation to unseen readers. Third, the GRU-based sequence model explored during development warrants further investigation with proper padded-batch training and a competition-metric early-stopping criterion.

\section{Conclusion}

We presented a multi-stage pipeline for predicting word-level Total Reading Time in Romanian texts, combining handcrafted linguistic features, causal surprisal from a Romanian GPT-2, masked language model features from Romanian BERT (following \citealt{hodivoianu2025predicting}), data augmentation, a two-stage LightGBM architecture, and session-aware rolling features. The pipeline is engineered around the core challenge of the task---generalising simultaneously to unseen readers and unseen texts---through leave-text-out feature computation, GroupKFold validation, and reader-profile augmentation. Our results confirm that LM-based predictability estimates are the most valuable contextual signal, consistent with prior findings in the reading-time literature.

\section{Limitations}

Our approach has several limitations. The rolling session features depend on the ordering of words within a participant's session; any corruption of this ordering would degrade predictions. The causal surprisal features require access to a Romanian autoregressive language model, which may be unavailable or slow on CPU-only hardware. The BERT feature extraction is computationally expensive and adds significant wall-clock time on machines without GPU acceleration. Finally, the leave-text-out aggregation assumes that the test texts are sufficiently similar in vocabulary to the training texts, which may not hold for highly domain-specific content.

\section{Ethical Statement}

This work uses anonymised eye-tracking data collected from consenting participants. No personally identifiable information is processed. The models and predictions are intended solely for research purposes within the NitroNLP hackathon context. We acknowledge that reading-time prediction models could in principle be misused to profile individual reading behaviour; we do not develop or release any such profiling capability.

\section{Acknowledgements}

We thank the organisers of the 5th Edition of the NitroNLP Hackathon for providing the dataset and evaluation infrastructure. The dataset originates from the MultiplEYE project (COST Action CA21131).

\begin{thebibliography}{10}

\bibitem{dumitrescu2020birth}
Stefan Dumitrescu, Andrei-Marius Avram, and Sampo Pyysalo. 2020.
\newblock The birth of {R}omanian {BERT}.
\newblock In \textit{Findings of the Association for Computational Linguistics: EMNLP 2020}, pages 4324--4328. Association for Computational Linguistics.

\bibitem{frazier1982making}
Lyn Frazier and Keith Rayner. 1982.
\newblock Making and correcting errors during sentence comprehension: Eye movements in the analysis of structurally ambiguous sentences.
\newblock \textit{Cognitive Psychology}, 14(2):178--210.

\bibitem{hodivoianu2025predicting}
Anamaria Hodivoianu, Oleksandra Kuvshynova, Filip Popovici, Adrian Luca, and Sergiu Nisioi. 2025.
\newblock Predicting total reading time using {R}omanian eye-tracking data.
\newblock In \textit{Proceedings of the First International Workshop on Gaze Data and NLP (Gaze4NLP) at RANLP 2025}, pages 71--75, Varna, Bulgaria.

\bibitem{hollenstein2021multilingual}
Nora Hollenstein, Federico Pirovano, Ce Zhang, Lena J{\"a}ger, and Lisa Beinborn. 2021.
\newblock Multilingual language models predict human reading behavior.
\newblock In \textit{Proceedings of the 2021 Conference of the North American Chapter of the Association for Computational Linguistics: Human Language Technologies}, pages 106--123. Association for Computational Linguistics.

\bibitem{dini2025human}
Luca Dini, Lucia Domenichelli, Dominique Brunato, and Felice Dell'Orletta. 2025.
\newblock From human reading to {NLM} understanding: Evaluating the role of eye-tracking data in encoder-based models.
\newblock In \textit{Proceedings of the 63rd Annual Meeting of the Association for Computational Linguistics (Volume 1: Long Papers)}, pages 17796--17813, Vienna, Austria. Association for Computational Linguistics.

\bibitem{jakobi2025multipleye}
Deborah Noemie Jakobi et al. 2025.
\newblock {MultiplEYE}: Creating a multilingual eye-tracking-while-reading corpus.
\newblock In \textit{Proceedings of the 2025 Symposium on Eye Tracking Research and Applications (ETRA '25)}. Association for Computing Machinery.

\bibitem{ke2017lightgbm}
Guolin Ke, Qi Meng, Thomas Finley, Taifeng Wang, Wei Chen, Weidong Ma, Qiwei Ye, and Tie-Yan Liu. 2017.
\newblock {LightGBM}: A highly efficient gradient boosting decision tree.
\newblock In \textit{Advances in Neural Information Processing Systems}, volume 30.

\bibitem{speer2022wordfreq}
Robyn Speer. 2022.
\newblock rspeer/wordfreq: v3.0.
\newblock \url{https://doi.org/10.5281/zenodo.7094679}.

\end{thebibliography}

\appendix

\section{Feature Summary}
\label{sec:appendix}

\begin{table}[h]
\centering
\small
\begin{tabular}{lll}
\toprule
\textbf{Group} & \textbf{Feature(s)} & \textbf{Source} \\
\midrule
Surface & word\_len, alpha\_len, n\_syllables & computed \\
Surface & is\_punct, is\_url, has\_digit & computed \\
Surface & is\_upper\_first, is\_all\_upper & computed \\
Surface & ends\_with\_punct & computed \\
Position & doc\_num, page\_num, word\_idx & word\_id parse \\
Position & pos\_in\_page, rel\_pos\_in\_page & computed \\
Frequency & log\_freq, prev/next log\_freq & corpus counts \\
Frequency & zipf\_freq, prev/next zipf\_freq & wordfreq \cite{speer2022wordfreq} \\
Per-word stats & tok\_mean, tok\_median, tok\_std & leave-text-out \\
Per-word stats & tok\_count, tok\_skip\_rate & leave-text-out \\
Genre & genre\_id, genre\_enc & leave-text-out \\
Context & prev/next surface features & shifted \\
Surprisal & surprisal, prev/next surprisal & RoGPT2 \\
BERT & mlm\_logp\_first, mlm\_logp\_sum & RoBERT \cite{hodivoianu2025predicting} \\
BERT & n\_bert\_subwords & RoBERT \\
BERT & bert\_pc\_0\ldots23 & PCA of embeddings \\
Rolling (pass 2) & session\_word\_idx, pass1\_pred & pass-1 OOF \\
Rolling (pass 2) & roll\_mean\_10/20/50 & rolling window \\
Rolling (pass 2) & roll\_skip\_10/20/50 & rolling window \\
\bottomrule
\end{tabular}
\caption{Complete feature list. ``leave-text-out'' means the statistic for a training row excludes the row's own text, preventing target leakage.}
\label{tab:features}
\end{table}

\end{document}
